% Section 13.8, translated from sv/chapters/13/exercises.tex.

\section{Exercises}
\label{c13:sec:uppgifter}

\begin{aside}
\textbf{How to work with the exercises.} Parts~1 and~2 are done during the lesson: first on your
own, a few minutes per exercise, then together. Part~3 is extra exercises for further practice
after the lesson; most of them can be done in your head or with pencil and paper.

\facitnotis{L13}
\end{aside}

% ------------------------------------------------------------------------------------------
\uppgiftsdel{Part 1}{Review exercises}

\uppgift{1.1}{Analysing a cubic function}
Consider the function
\[
  f(x) = -0.5x^3 + 1.5x^2 + 4.5x - 3.
\]
\begin{parts}
  \item Differentiate the function $f(x)$ and find the expression for $f'(x)$.
  \item Find where the function is stationary, that is, solve the equation $f'(x) = 0$.
  \item Use the second derivative to decide whether each stationary point is a maximum or a
    minimum.
  \item Calculate the maximum and minimum values of the function, that is, the function values at
    the stationary points.
  \item Plot the graph of $f(x)$ in GeoGebra, \url{https://www.geogebra.org/graphing?lang=en}, and
    check that your answers are right.
\end{parts}

\uppgift{1.2}{Analysing the voltage across a capacitor}
A capacitor in an RC circuit is charged and discharged repeatedly. The voltage across the
capacitor can be described by the function
\[
  u(t) = -0.2t^3 + 1.2t^2 - 1.4625t + 3,
\]
where $u(t)$ is the voltage across the capacitor in V and $t$ the time in seconds,
$0 \leq t \leq 6$.
\begin{parts}
  \item Differentiate the function and find the expression for the rate of change of the voltage
    $u'(t)$.
  \item Find the time (or times) at which the voltage is stationary, that is, when
    $u'(t) = 0$.
  \item Use the second derivative to decide whether the points are maximum or minimum points.
  \item Calculate the maximum and minimum values of the voltage, that is, the voltage at the
    stationary points.
  \item Plot the graph of $u(t)$ in GeoGebra, \url{https://www.geogebra.org/graphing?lang=en}, and
    check that your answers are right.
\end{parts}

% ------------------------------------------------------------------------------------------
\uppgiftsdel{Part 2}{New material}

\uppgift{2.1}{Current through an inductor}
The current through an inductor varies over time according to the function
\[
  i(t) = 4\sin(50t),
\]
where $i(t)$ is the current through the inductor in A and $t$ the time in seconds.
\begin{parts}
  \item Differentiate the function and find the expression for $i'(t)$.
  \item Find the time (or times) at which the current is stationary, that is, when
    $i'(t) = 0$.
  \item Use the second derivative to decide whether the points are maximum or minimum points.
  \item Calculate the greatest and least values of the current. Consider whether this matches
    what you expected.
\end{parts}

\uppgift{2.2}{Discharging a capacitor}
The voltage across a capacitor that is discharging is given by the function
\[
  u(t) = 12e^{-0.4t},
\]
where $u(t)$ is the voltage across the capacitor in V and $t$ the time in seconds.
\begin{parts}
  \item Differentiate the function and find the expression for $u'(t)$.
  \item Calculate the rate of change of the voltage at $t = 2\,\text{s}$.
  \item State whether the rate of change of the voltage is increasing or decreasing at this
    point, that is, find $u''(2)$.
\end{parts}

\uppgift{2.3}{Gain in dB}
The logarithmic gain of an amplifier system is described by the function
\[
  G(x) = 10\log_{10}(2x + 1),
\]
where $G(x)$ is the gain in dB and $x$ the input level (a dimensionless factor).
\begin{parts}
  \item Calculate the value of $x$ for which the gain is $G(x) = 2\,\text{dB}$.
  \item Calculate the rate of change of the gain $G'(x)$ at this point.
\end{parts}

\begin{aside}
\note[Note] The rate of change $G'(x)$ states how fast the gain changes when the input $x$
changes.
\end{aside}

% ------------------------------------------------------------------------------------------
\uppgiftsdel{Part 3}{Extra exercises}
The exercises below are extra practice, best done on your own after the lesson.

\uppgift{3.1}{Differentiating trigonometric functions}
Differentiate.
\begin{delar}(3)
  \task $f(x) = \sin(3x)$
  \task $f(x) = \cos(5x)$
  \task $f(x) = 4\sin(x)$
  \task $f(x) = -2\cos(2x)$
  \task $f(x) = \sin(x) + \cos(x)$
\end{delar}

\uppgift{3.2}{Differentiating exponential functions}
Differentiate.
\begin{delar}(3)
  \task $f(x) = e^{5x}$
  \task $f(x) = 3e^{-2x}$
  \task $f(x) = e^x + x^2$
  \task $f(x) = 2^x$
  \task $f(x) = 10e^{-x/4}$
\end{delar}

\uppgift{3.3}{Differentiating logarithmic functions}
Differentiate.
\begin{delar}(3)
  \task $f(x) = \ln x$
  \task $f(x) = \ln(7x)$
  \task $f(x) = 3\ln x$
  \task $f(x) = \log_{10} x$
  \task $f(x) = \ln x + e^x$
\end{delar}

\uppgift{3.4}{Mixed differentiation}
Differentiate.
\begin{delar}(2)
  \task $f(x) = 2e^{3x} - \sin(4x)$
  \task $f(x) = x^2 + \ln(2x)$
  \task $f(x) = 5\cos(x) - 3e^{-x}$
  \task $f(x) = 4\sin(2x) + 4\cos(2x)$
\end{delar}

\uppgift{3.5}{The derivative at a point}
Calculate the derivative at the given point.
\begin{delar}(3)
  \task $f(x) = e^{2x}$ at $x = 0$
  \task $f(x) = \sin x$ at $x = \dfrac{\pi}{2}$
  \task $f(x) = \ln x$ at $x = 2$
  \task $f(x) = \cos(3x)$ at $x = 0$
  \task $f(x) = 3e^{-x}$ at $x = 1$
\end{delar}

\uppgift{3.6}{Current through a capacitor}
The voltage across a capacitor with $C = 47\,\mu\text{F}$ is given by $u(t) = 10\sin(100\pi t)$
volts. The current is given by $i_C(t) = C\dfrac{du}{dt}$.
\begin{parts}
  \item Find $u'(t)$.
  \item Find an expression for $i_C(t)$.
  \item What is the amplitude of the current?
  \item At what time is the current greatest for the first time?
\end{parts}

\uppgift{3.7}{Voltage across an inductor}
The current through an inductor with $L = 25\,\text{mH}$ is given by $i(t) = 2\sin(200t)$
amperes. The voltage is given by $u_L(t) = L\dfrac{di}{dt}$.
\begin{parts}
  \item Find $i'(t)$.
  \item Find an expression for $u_L(t)$.
  \item What is the amplitude of the voltage?
  \item What is the phase difference between the voltage and the current?
\end{parts}

\uppgift{3.8}{Discharging a capacitor}
The voltage is given by $u(t) = 20e^{-t/0.5}$ volts.
\begin{parts}
  \item Find $u'(t)$.
  \item Calculate $u'(0)$.
  \item Calculate $u'(0.5)$.
  \item What does the sign of the derivative tell you?
  \item Show that $u'(t) = -\dfrac{u(t)}{\tau}$ with $\tau = 0.5\,\text{s}$.
\end{parts}

\uppgift{3.9}{Stationary point of an exponential function}
The function $f(x) = xe^{-2x}$ is given.
\begin{parts}
  \item Find $f'(x)$ with the product rule and factor out $e^{-2x}$.
  \item Solve $f'(x) = 0$.
  \item Use $f''(x)$ to decide whether the point is a maximum or a minimum.
  \item Calculate the extreme value of the function.
\end{parts}

\uppgift{3.10}{Maximum power over time}
The power in a component is given by $p(t) = 100te^{-t}$ watts for $t \geq 0$.
\begin{parts}
  \item Find $p'(t)$.
  \item At what time is the power stationary?
  \item Use $p''(t)$ to decide whether it is a maximum or a minimum.
  \item Calculate the maximum power.
\end{parts}

\uppgift{3.11}{The derivative of a dB function}
The gain is given by $G(x) = 20\log_{10}(x)$ dB.
\begin{delar}(3)
  \task Find $G'(x)$.
  \task Calculate $G'(1)$.
  \task Calculate $G'(10)$.
  \task* Interpret the difference between the answers to \textbf{b)} and \textbf{c)}.
\end{delar}

\needspace{10\baselineskip}
\uppgift{3.12}{Tangent to an exponential curve}
The function $f(x) = e^{-x}$ is given. Consider the point where $x_0 = 0$.
\begin{parts}
  \item Calculate $f(0)$.
  \item Calculate $f'(0)$.
  \item Find the equation of the tangent.
  \item Where does the tangent cross the $x$-axis? Compare with the time constant $\tau$ of a
    discharge.
\end{parts}

\uppgift{3.13}{The derivative of an exponential discharge}
The voltage is given by $u(t) = U_0e^{-t/\tau}$.
\begin{parts}
  \item Show that $u'(t) = -\dfrac{u(t)}{\tau}$.
  \item What does this relationship mean physically?
  \item Calculate $u'(0)$ and $u'(\tau)$ when $U_0 = 10\,\text{V}$ and $\tau = 2\,\text{s}$.
  \item What fraction of the initial voltage remains after the time $\tau$?
\end{parts}

\uppgift{3.14}{True or false}
Decide whether each statement is true or false, and justify your answer briefly.
\begin{parts}
  \item The derivative of $\sin x$ is $-\cos x$.
  \item The derivative of $e^{3x}$ is $e^{3x}$.
  \item The derivative of $\ln(5x)$ is $\dfrac{1}{5x}$.
  \item The derivative of $\cos(2x)$ is $-2\sin(2x)$.
  \item $e^x$ is its own derivative.
  \item The derivative of a constant times a function is the constant times the derivative of the
    function.
\end{parts}

\uppgift{3.15}{Find the error}
Each line contains a common error. Explain the error and give the correct answer.
\begin{delar}(2)
  \task $\dfrac{d}{dx}\sin(3x) = \cos(3x)$
  \task $\dfrac{d}{dx}e^{-2x} = -2e^{-2x-1}$
  \task $\dfrac{d}{dx}\ln(3x) = \dfrac{1}{3x}$
  \task $\dfrac{d}{dx}4\cos(x) = 4\sin(x)$
  \task $\dfrac{d}{dx}2^x = x \cdot 2^{x-1}$
\end{delar}
